\subsubsection{Module Overview}

In the data science lifecycle, data understanding, exploratory data analysis (EDA), and data preparation represent major stages that involve a large amount of repetitive work. In this module, the aim is to reduce this repetitive effort by introducing a rule-based system that automates these processes.

The module is designed to work based on the type of data, whether tabular, text, or image. For each data type, the corresponding data understanding and preprocessing steps are applied to ensure that the data is properly prepared and suitable for input into machine learning and deep learning models.

The system also generates a final report that provides the user with information about the dataset and the applied preprocessing steps. For example, in tabular data, this includes details such as missing values, number of duplicates, dataset size, and column types. In addition, it describes the preprocessing steps that were applied to the data and presents the final results after processing.
\subsubsection{Supported Dataset Types}
\begin{itemize}
  \item Tabular Dataset
  \item Text Dataset
  \item Image Dataset
\end{itemize}
\subsubsection{Supported Tasks }

\begin{itemize}
  \item Classification and Regression for tabular dataset
  \item Classification for text dataset
  \item Classification and Segmentation for image dataset

\end{itemize}

\subsubsection{Tabular Dataset}

The meaning of The tabular dataset in this module is structured data organized in rows and columns, where features can be either categorical or numerical.
\vspace{1cm}

\Needspace{5\baselineskip}
\noindent\textbf{Training Preprocessing Steps}\par
\begin{enumerate}
  \item Data Input and User Configuration Initialization.
  \item Get data overview.
  \item Drop duplicates.
  \item Split dataset.
  \item Get training data information.
  \item Specify columns to drop.
  \item Drop selected columns.
  \item Detect the type of each column based on the training data.
  \item Generate visualizations of the training data based on column types and task type.
  \item Apply feature preprocessing techniques for each column based on its detected type.
  \item Apply target preprocessing based on the selected task.
  \item Save all files required for preprocessing the test data.
  \item Generate Report if user enable report generation.
  \item Return \texttt{X\_train}, \texttt{y\_train}, \texttt{X\_val}, \texttt{y\_val}.
\end{enumerate}
\vspace{1cm}
\Needspace{4\baselineskip}
\noindent\textbf{Data Input and User Configuration Initialization:}\par
The preprocessing pipeline is initialized through the class constructor. 
During initialization, all user-defined parameters are assigned and validated to ensure consistency and correctness before executing any preprocessing steps. 
These validations help prevent invalid configurations such as incorrect problem types, wrong target column name, or invalid parameter values, which could otherwise lead to runtime errors or incorrect model behavior.
\begin{table}[!htbp]
\centering

\caption{Classical Training Data Preprocessing Parameters - Description}
\label{tab:params_description}

\resizebox{\textwidth}{!}{%
\begin{tabular}{ll}
\toprule
\textbf{Parameter} & \textbf{Description} \\
\midrule

df &
The dataframe of the dataset. \\

target\_col &
The name of the target column in the dataset. \\

problem\_type &
Classification or Regression task type. \\

folder\_path &
Path to store reports, charts, and serialized objects. \\

val\_size &
Percentage of data used for validation split. \\

random\_state &
To ensure reproducible data splitting. \\

cardinality\_threshold &
Threshold for separating low and high cardinality categorical features. \\

max\_unique\_ordinal &
Maximum number of unique values to treat a column as ordinal. \\

missing\_threshold &
Threshold for dropping columns based on missing values ratio. \\


columns\_need\_to\_drop &
List of user-specified columns to remove. \\

is\_report &
Enables or disables preprocessing report generation. \\

\bottomrule
\end{tabular}%
}
\end{table}

\begin{table}[!htbp]
\centering

\caption{Classical Training Data Preprocessing Parameters - Constraints}
\label{tab:params_constraints}

\resizebox{\textwidth}{!}{%
\begin{tabular}{ll}
\toprule
\textbf{Parameter} & \textbf{Allowed Values / Constraints} \\
\midrule

df &
pandas DataFrame (must include required columns) \\

target\_col &
String; must exist in dataframe columns \\

problem\_type &
``classification'' or ``regression'' \\

folder\_path &
String path  \\

val\_size &
Float: $0 < val\_size \leq 0.5$ \\

random\_state &
Integer or None \\

cardinality\_threshold &
Integer $\geq 0$ \\

max\_unique\_ordinal &
Integer $\geq 0$ \\

missing\_threshold &
Float: $0 \leq x \leq 1$ \\

columns\_need\_to\_drop &
List of strings \\

is\_report &
Boolean (True / False) \\

\bottomrule
\end{tabular}%
}
\end{table}
\Needspace{4\baselineskip}
\vspace{1cm}
\noindent\textbf{Data Overview:}\par

This step is responsible for generating a comprehensive statistical and structural summary of the dataset before preprocessing begins, helping to create a detailed report.
\begin{enumerate}
    \item \textbf{Capture initial data snapshot:} The first five rows of the original dataset are stored to preserve the original structure before any modifications.

    \item \textbf{Standardize text data:} All string-based values in the dataset are converted to lowercase, ensuring consistency in categorical values.

    \item \textbf{Capture transformed snapshot:} After applying lowercase transformation, another preview of the dataset is stored for comparison purposes.

    \item \textbf{Generate dataset structure information:} to capture column types, non-null counts, and memory usage.

    \item \textbf{Separate numerical and categorical features:} Columns are divided based on data type into numerical and categorical subsets using \texttt{select\_dtypes()}.

    \item \textbf{Compute descriptive statistics:}
    \begin{itemize}
        \item Numerical features: mean, standard deviation, min, max, quartiles.
        \item Categorical features: frequency, unique values, top category.
    \end{itemize}

    \item \textbf{Compute missing values:} The number of missing values per column is calculated.

    \item \textbf{Detect duplicate rows:}
    \begin{itemize}
        \item Total number of duplicate rows is computed.
        \item Percentage of duplicated rows is calculated relative to dataset size.
    \end{itemize}

    \item \textbf{Store dataset shape:} The dataset dimensions (rows and columns) are recorded.

    \item \textbf{Save all results into report structure:} All computed statistics, snapshots, and metadata are stored for later reporting and visualization.
\end{enumerate}
\begin{lstlisting}[style=algostyle,caption={Data Overview Generation},label={alg:data_overview}]
Input: Dataframe D
Output: Stored data overview in report object

1:  Extract the first rows of D as an initial preview
2:  Convert string values in object-type columns to lowercase
3:  Extract dataframe structure information using the info function
4:  Split D into numerical features D_num and categorical features D_cat
5:  if D_num is not empty then
6:      compute numerical statistics: mean, standard deviation, min, max, quartiles
7:  end if
8:  if D_cat is not empty then
9:      compute categorical statistics: unique values, frequencies, top category
10: end if
11: Compute missing values for each column
12: Compute duplicate row count and duplicate percentage
13: Store dataframe shape: number of rows and columns
14: Save head, info, statistics, missing values, duplicates, and shape in report_data
15: Return the updated report object
\end{lstlisting}
\Needspace{4\baselineskip}
\noindent\textbf{Drop duplicates:}\par

Duplicate rows are removed from the dataset to prevent the model from overfitting to repeated examples and to ensure a more reliable evaluation process.

After performing this step, the system reports key statistics to the user, including the total number of duplicate rows detected, the percentage of duplicated data, and the resulting dataset shape after removal. This allows the user to verify that the dataset no longer contains redundant entries.
\vspace{1cm}
\Needspace{4\baselineskip}
\noindent\textbf{Split data:}\par
Split the data into \texttt{X\_train}, \texttt{X\_val},
\texttt{y\_train}, and \texttt{y\_val} using the
\texttt{train\_test\_split} function from Scikit-learn. The split follows the
configured \texttt{test\_size} and \texttt{random\_state} values.
\vspace{1cm}
\Needspace{4\baselineskip}
\noindent\textbf{Get training data information:}\par
 For each column, calculate the required information,
such as the percentage of missing values,number of unique values, the mean and median for
numerical columns, and the mode value
\Needspace{4\baselineskip}
\vspace{1cm}
\noindent\textbf{Specify Columns to Drop:}\par
A column is dropped if:
\begin{itemize}
    \item It is included in the user-provided \texttt{columns\_need\_to\_drop} parameter.
    \item The percentage of missing values is greater than the user-provided \texttt{missing\_threshold}.
    \item The column is a constant value.
\end{itemize}
\Needspace{4\baselineskip}
\noindent\textbf{Drop Selected Columns:}\par
The specified columns are removed from both \texttt{X\_train} and \texttt{X\_val}.The dropped column names are serialized and saved as a pickle file to ensure consistency when applying the same preprocessing steps on future data.
\Needspace{4\baselineskip}
\vspace{1cm}
\noindent\textbf{Detect the type of each column based on the training data:}\par
Each feature is automatically assigned a column type. The detected type determines the preprocessing techniques that will be applied later in the pipeline.
\begin{table}[!htbp]
\centering
\caption{Detected column types.}
\begin{tabular}{p{0.25\linewidth}p{0.70\linewidth}}
\toprule
\textbf{Column type} & \textbf{Detection rule}\\
\midrule
\textbf{Binary} & Columns containing exactly two unique values or boolean values.\\
\textbf{Ordinal} & Integer columns with a number of unique values less than or equal to \texttt{max\_unique\_ordinal}.\\
\textbf{Numerical} & Integer columns with a number of unique values greater than \texttt{max\_unique\_ordinal}.\\
\textbf{Continuous} & Floating-point columns.\\
\textbf{Low-Cardinality Categorical} & Object columns with a number of unique values less than or equal to \texttt{cardinality\_threshold}.\\
\textbf{High-Cardinality Categorical} & Object columns with a number of unique values greater than \texttt{cardinality\_threshold}.\\
\textbf{Other} & Columns with unsupported data types that are excluded from the preprocessing pipeline.\\
\bottomrule
\end{tabular}
\end{table}
\begin{itemize}
    \item \textbf{Binary}: Columns containing exactly two unique values or boolean values.  
    Examples: ``yes/no'', ``1/0'', ``True/False''.

    \item \textbf{Ordinal}: Integer columns with a small number of unique values less than or equal \texttt{max\_ordinal\_unique}, where values have a meaningful order.  
    Example: $1, 2, 3, 4$ (ratings, levels).

    \item \textbf{Numerical}: Integer columns with a number of unique values greater than \texttt{max\_ordinal\_unique}, typically representing continuous or high-variance numeric data.  
    Examples: age, number of transactions.

    \item \textbf{Continuous}: Floating-point columns representing real-valued measurements.  
    Examples: height = 175.5, price = 19.99, temperature = 36.6.

    \item \textbf{Low-Cardinality Categorical}: Object (string) columns with a number of unique values less than or equal to \texttt{cardinality\_threshold}.  
    Examples: product categories.

    \item \textbf{High-Cardinality Categorical}: Object (string) columns with a number of unique values greater than \texttt{cardinality\_threshold}.  
    Examples: product names, cities.

    \item \textbf{Other}: Columns with unsupported or non-standard data types that are excluded from the preprocessing pipeline.  
  
\end{itemize}
\vspace{1cm}
\begin{lstlisting}[style=algostyle,caption={Column Type Detection and Preprocessing Assignment}]
Input: Training features X_train
Output: Column groups and assigned preprocessing rules

1:  for each column in X_train do
2:      if column is binary then
3:          assign type Binary
4:          assign most-frequent imputation and ordinal encoding
5:      else if column is integer then
6:          if unique values <= maximum ordinal threshold then
7:              assign type Ordinal
8:              assign most-frequent imputation and robust scaling
9:          else
10:             assign type Numerical
11:             assign median imputation and robust scaling
12:         end if
13:     else if column is floating-point then
14:         assign type Continuous
15:         assign median imputation and robust scaling
16:     else if column is categorical then
17:         if cardinality <= cardinality threshold then
18:             assign type Low-Cardinality Categorical
19:             assign most-frequent imputation and binary encoding
20:         else
21:             assign type High-Cardinality Categorical
22:             assign most-frequent imputation and frequency encoding
23:         end if
24:     else
25:         mark column as unsupported and schedule it for removal
26:     end if
27:     store column type and preprocessing steps in the report
28: end for
29: Return all categorized column groups
\end{lstlisting}

\Needspace{4\baselineskip}
\noindent\textbf{Generate visualizations of the training data based on column types and task type:}\par
When the \texttt{is\_report} flag is enabled, the tabular preprocessing report produces a set of visualizations corresponding to each detected feature type and task type.
\begin{itemize}
  \item \textbf{Binary / categorical features with classification targets:} The report produces these graphs
    \begin{itemize}
      \item a null-vs-non-null pie chart,
      \item a category count plot ,
      \item a category vs target count plot .
    \end{itemize}
  \item \textbf{Binary / categorical features with regression targets:} The report produces these graphs
    \begin{itemize}
      \item a null-vs-non-null pie chart,
      \item a category count plot ,
      \item a box plot of the continuous target for each category.
    \end{itemize}
  \item \textbf{Ordinal features with classification targets:} The report produces these graphs
    \begin{itemize}
      \item a null-vs-non-null pie chart,
      \item an ordered count plot,
      \item an ordered count plot split by class label.
    \end{itemize}
  \item \textbf{Ordinal features with regression targets:} The report  produces these graphs
    \begin{itemize}
      \item a null-vs-non-null pie chart,
      \item an ordered count plot,
      \item a box plot of the regression target for each ordinal level.
    \end{itemize}
  \item \textbf{Numerical / continuous features with classification targets:} The report  produces these graphs
    \begin{itemize}
      \item a null-vs-non-null pie chart,
      \item a histogram with a kernel density estimate,
      \item a box plot of the feature grouped by class label.
    \end{itemize}
  \item \textbf{Numerical / continuous features with regression targets:} The report  produces these graphs
    \begin{itemize}
      \item a null-vs-non-null pie chart,
      \item a histogram with a kernel density estimate,
      \item a scatter plot of feature value against the regression target.
    \end{itemize}
  \item \textbf{Target variable for classification:} The report produces these graphs
    \begin{itemize}
      \item a null-vs-non-null pie chart,
      \item a target class frequency count plot.
    \end{itemize}
  \item \textbf{Target variable for regression:} The report produces these graphs
    \begin{itemize}
      \item a null-vs-non-null pie chart,
      \item a histogram with a kernel density estimate,
      \item a box plot of the target distribution.
    \end{itemize}
  \item \textbf{Numeric correlation matrix:} The report generates a heatmap of pairwise correlations across all numeric features.
\end{itemize}
\Needspace{2\baselineskip}
\textbf{Interpretation of visualization types:}
\begin{itemize}
  \item \textbf{Null vs non-null pie chart:} shows the proportion of missing values in a feature, helping detect data quality issues.

  \item \textbf{Category count plot:} shows the frequency distribution of categorical values and reveals imbalance or rare categories.

  \item \textbf{Category vs target count plot:} shows how each category relates to the target variable,helps identify whether certain categories are strongly associated with specific target outcomes.

  \item \textbf{Ordered count plot:} shows the frequency of ordinal values while preserving their order, helping visualize how samples are distributed across ordered levels.

  \item \textbf{Ordered count plot split by class:} shows how different classes are distributed across ordered levels (e.g., 1 to 4), helping identify whether certain classes tend to appear more in specific ordinal levels.

  \item \textbf{Box plot (category/ordinal vs target):} shows the distribution of a continuous target across groups, highlighting differences, variance, and outliers.

  \item \textbf{Histogram + KDE:} shows the distribution shape of a feature or target, including skewness, modality, and spread.

  \item \textbf{Scatter plot (feature vs target):} shows relationships or trends between numerical features and regression targets, including linear/non-linear patterns.

  \item \textbf{Box plot grouped by class label:} shows how numerical features vary across classes and helps detect outliers.

  \item \textbf{Correlation heatmap:} shows pairwise linear relationships between numerical features, helping detect multicollinearity and redundancy.

  \item \textbf{Target frequency plot (classification):} shows class imbalance, which is critical for model selection and evaluation strategy.

  \item \textbf{Target histogram (regression):} shows distribution of target values, helping detect skewness and outliers.
\end{itemize}
\Needspace{4\baselineskip}
\noindent\textbf{Apply feature preprocessing techniques for each column based on its detected type:}\par

\begin{table}[!htbp]
\centering
\caption{Feature preprocessing applied by column type.}
\begin{tabular}{p{0.30\linewidth}p{0.65\linewidth}}
\toprule
\textbf{Column type} & \textbf{Applied preprocessing}\\
\midrule
\textbf{Numerical / continuous} & Imputation with \texttt{SimpleImputer(strategy="median")}, followed by \texttt{RobustScaler()}.\\
\textbf{Binary} & Most-frequent imputation followed by ordinal encoding with unknown values mapped to -1.\\
\textbf{Low-cardinality categorical} & Most-frequent imputation followed by binary encoding.\\
\textbf{Frequency-encoded categorical} & Most-frequent imputation followed by a frequency encoding transform.\\
\textbf{Ordinal} & Most-frequent imputation followed by robust scaling.\\
\bottomrule
\end{tabular}
\end{table}
\begin{table}[!htbp]
\centering
\caption{Justification of preprocessing techniques by feature type.}
\begin{tabular}{p{0.30\linewidth}p{0.65\linewidth}}
\toprule
\textbf{Preprocessing choice} & \textbf{Justification} \\
\midrule

Median imputation &
Robust to outliers compared to mean imputation, making it suitable for skewed numerical distributions. \\

Robust scaling &
Reduces the influence of outliers by scaling using the interquartile range instead of mean and variance. \\

Most-frequent imputation &
Preserves the natural distribution of categorical variables by replacing missing values with the dominant category. \\

Binary encoding &
Efficient representation for categorical variables, reducing dimensionality from $O(k)$ in one-hot encoding to $O(\log k)$ while preserving information. \\

Frequency encoding &
Captures category importance based on occurrence frequency, reducing sparsity and preserving distributional signals. \\

Ordinal handling &
Preserves inherent ordering information in the feature (e.g., 1 < 2 < 3 < 4), unlike nominal encodings. \\

Unknown category mapping (-1) &
Ensures robustness during inference by handling unseen categories without breaking the encoding schema. \\

\bottomrule
\end{tabular}
\end{table}
\Needspace{4\baselineskip}
\noindent\textbf{Apply target preprocessing based on the selected task:}\par
Target preprocessing changes based on problem type (classification / regression):

\begin{itemize}
  \item \textbf{Classification:} Fill missing target values with the training-mode label, then encode the labels with \texttt{LabelEncoder}.
  \item \textbf{Regression:} Fill missing target values with the training median, then standard scale the target using \texttt{StandardScaler}. 
\end{itemize}
\begin{lstlisting}[style=algostyle,caption={Target Variable Preprocessing}]
Input: Target y_train, y_val and problem type
Output: Processed targets and saved target metadata

1:  if problem type is Classification then
2:      replace missing target values using the most frequent target value
3:      fit a LabelEncoder on the training target
4:      transform training and validation targets
5:      save the fitted LabelEncoder
6:      record classification target preprocessing in the report
7:  else if problem type is Regression then
8:      replace missing target values using the median target value
9:      fit a StandardScaler on the training target
10:     scale training and validation targets
11:     save the fitted StandardScaler
12:     record regression target preprocessing in the report
13: end if
14: Store processed target samples in the report
15: Save target metadata for inference
16: Return processed training and validation targets
\end{lstlisting}

\vspace{1cm}
\Needspace{4\baselineskip}
\noindent\textbf{Save all files required for preprocessing the test data:}\par
The preprocessing pipeline persists the objects required for later inference and reuse:

\begin{table}[!htbp]
\centering
\caption{Saved preprocessing artifacts for test-time reuse.}
\begin{tabular}{p{0.30\linewidth}p{0.65\linewidth}}
\toprule
\textbf{Artifact} & \textbf{Purpose}\\
\midrule
\texttt{training\_preprocessor.}\newline\texttt{pkl} & The fitted feature preprocessor.\\
\texttt{feature\_names.pkl} & The output feature names produced by the feature preprocessor.\\
\texttt{target\_preprocessor.}\newline\texttt{pkl} & The fitted target encoder or scaler.\\
\texttt{target\_info.pkl} & Information about the target, including target name, problem type, mode, and median.\\
\texttt{bool\_type\_col.pkl} & Boolean columns that are converted to integer values during preprocessing.\\
\texttt{data\_columns.pkl} & The column names of the input DataFrame before preprocessing.\\
\texttt{col\_drop.pkl} & The names of the dropped columns.\\
\bottomrule
\end{tabular}
\end{table}
\vspace{1cm}
\Needspace{4\baselineskip}
\noindent\textbf{Generate Report (Optional):}\par
When \texttt{is\_report} is enabled, the \texttt{TabularPreprocessingReport} class generates an HTML report from the collected \texttt{report\_data}. The report documents the complete preprocessing workflow, including data overview statistics, missing values, duplicate analysis, train-validation split details, removed columns, generated graphs, preprocessing decisions for each feature, and previews of the transformed training and validation datasets. This allows users to inspect and understand all preprocessing steps applied before model training.
\vspace{1cm}
\Needspace{5\baselineskip}
\noindent\textbf{Testing Preprocessing Steps}\par
The testing preprocessing workflow is designed to reuse the exact training preprocessing artifacts and transformations.

\begin{enumerate}
  \item Verify that the required artifact files exist:
        target preprocessor, training preprocessor,
        \texttt{data\_columns.pkl}, \texttt{col\_drop.pkl},
        \texttt{target\_info.pkl}, and \texttt{bool\_type\_col.pkl}.
  \item Load the saved dataset columns, dropped columns, target information,
        boolean-column list, training preprocessor, and target preprocessor.
  \item Confirm \texttt{target\_info} contains a valid \texttt{col\_name} and \texttt{problem\_type} (classification or regression).
  \item Confirm \texttt{target\_preprocessor} and \texttt{training\_preprocessor} is not none.
  \item Ensure that the test dataset column order matches the saved training columns. If the target column is missing, testing proceeds in \texttt{features\_only} mode.
  \item Normalize string values in the test dataset by lowercasing them, matching the same text normalization used during training preprocessing.
  \item Separate the test DataFrame into \texttt{X\_test} and \texttt{y\_test} when \texttt{features\_only} mode is false.
\item Convert the columns listed in \texttt{bool\_type\_col.pkl} to integer data type.  \item Drop the same columns identified during training using \texttt{col\_drop.pkl}.
  \item Apply the saved \texttt{training\_preprocessor.pkl} transformer to \texttt{X\_test}.
  \item Apply target preprocessing when  when \texttt{features\_only} mode is false:
    \begin{itemize}
      \item For classification, fill missing labels with the training-mode value, then apply the saved label encoder.
      \item For regression, fill missing values with the training median, then apply the saved standard scaler.
    \end{itemize}
\item Return the preprocessed test features and labels when \texttt{features\_only} is set to \texttt{False}; otherwise, return the preprocessed test features and \texttt{None} for the labels.
\end{enumerate}

This ensures the test pipeline uses the same feature transformation and target encoding logic as training.
\begin{table}[!htbp]
\centering

\caption{Classical Testing Data Preprocessing Parameters}
\label{tab:classicaltestingparams}

\resizebox{\textwidth}{!}{%
\begin{tabular}{p{0.30\linewidth}p{0.65\linewidth}}

\toprule
\textbf{Parameter} & \textbf{Description} \\
\midrule

df & The dataframe of the test dataset It can cotain target column also can not in this case enable feature\_only mode. \\

folder\_path & Path to load saved preprocessing artifacts. \\

\bottomrule
\end{tabular}%
}
\end{table}

\subsubsection{Text Dataset}

The text dataset consists of unstructured data containing a single text column used for natural language processing tasks.

\vspace{1cm}
\Needspace{5\baselineskip}
\noindent\textbf{Training Preprocessing Steps}\par

\begin{enumerate}
  \item Constructor Initialization and Input Validation.
  \item Initialize NLP dependencies.
  \item Generate data overview and split dataset.
  \item Generate optional graphs.
  \item Handle missing text values.
  \item Apply text preprocessing and tokenization.
  \item Transform text using TF-IDF vectorization.
  \item Encode target labels.
  \item Save preprocessing artifacts.
  \item Generate preprocessing report (optional).
  \item Return processed datasets.
\end{enumerate}
\vspace{1cm}
\vspace{0.5cm}
\noindent\textbf{Constructor Initialization and Input Validation:}\par

The constructor initializes all required parameters. During initialization, the method performs  validation checks to ensure that the text column is correctly defined, is different from the target column, and has the appropriate data type. It also verifies that the target variable is suitable for classification tasks. Additionally, unnecessary columns are removed, required NLP resources are downloaded, and all configuration parameters are stored for later use.

\begin{table}[!htbp]
\centering
\caption{Text Training Data Preprocessing Parameters }
\label{tab:text_constructor_params}

\begin{tabular}{p{0.30\linewidth}p{0.65\linewidth}}
\toprule
\textbf{Parameter} & \textbf{Description} \\
\midrule

\texttt{df} &
Input pandas DataFrame containing the dataset. \\

\texttt{target\_col} &
Name of the target column used for classification. \\

\texttt{text\_col} &
Name of the column containing raw text data for NLP processing. \\

\texttt{folder\_path} &
Directory path used to store preprocessing outputs and reports. \\

\texttt{val\_size} &
Fraction of data used for validation split (default: 0.2). \\

\texttt{random\_state} &
Random seed for reproducible train-validation splitting (default: 42). \\

\texttt{tfidf\_max\_features} &
Maximum vocabulary size for TF-IDF vectorization (default: 500). \\

\texttt{tfidf\_ngram} &
N-gram range for TF-IDF vectorizer (default: (1,2)). \\

\texttt{tfidf\_max\_df} &
Maximum document frequency threshold for TF-IDF (default: 0.85). \\

\texttt{tfidf\_min\_df} &
Minimum document frequency threshold for TF-IDF (default: 3). \\

\texttt{is\_report} &
Boolean flag to enable or disable report generation. \\

\bottomrule
\end{tabular}
\end{table}
\vspace{1cm}

\Needspace{4\baselineskip}
\noindent\textbf{Initialize NLP dependencies:}\par
Required Natural Language Processing resources such as tokenizers and stopword lists are downloaded and initialized to enable text preprocessing operations.
\vspace{1cm}
\Needspace{4\baselineskip}
\noindent\textbf{Generate data overview and split dataset:}\par
A statistical overview of the dataset is generated, including missing values, duplicates, and descriptive statistics. After that, the dataset is split into training and validation sets using the configured split ratio.
\vspace{1cm}
\Needspace{4\baselineskip}
\noindent\textbf{Generate optional graphs:}\par
If enabled, the system generates visualizations such as text distribution plots and target distribution plots to better understand dataset characteristics.
\vspace{0.5cm}

\begin{table}[!htbp]
\centering
\caption{Optional Graph Generation in Text Preprocessing}
\begin{tabular}{p{0.35\linewidth}p{0.60\linewidth}}
\toprule
\textbf{Graph Type} & \textbf{Purpose} \\
\midrule

Null vs Non-Null Distribution &
Shows missing value ratio in text and target columns. \\

Target Distribution Plot &
Displays class frequency distribution for classification tasks. \\

Top-K Word Frequency Plot &
Shows the most frequent 7 words in the corpus.\\

\bottomrule
\end{tabular}
\end{table}
\vspace{1cm}
\Needspace{4\baselineskip}
\noindent\textbf{Handle missing text values:}\par
All missing values in the text column are replaced with empty strings.
\vspace{1cm}
\Needspace{4\baselineskip}
\noindent\textbf{Apply text preprocessing and tokenization:}\par
Raw text is cleaned and transformed using a custom tokenizer. This includes lowercasing, removing noise characters, handling contractions, removing stopwords, and applying stemming to normalize tokens.
\begin{lstlisting}[style=algostyle,caption={Custom Text Tokenizer},label={alg:custom_text_tokenizer}]
Input: Raw text string
Output: List of cleaned and stemmed tokens

1:  Initialize English stopwords and remove {"not", "no"} from stopword list
2:  Initialize PorterStemmer
3:  Convert text to lowercase and strip whitespace
4:  Normalize special apostrophes to signle qoute
5:  Replace contractions:
        "can't" --> "can not"
        "won't" --> "will not"
        "shan't" --> "shall not"
6:  Expand "n't" patterns into " not "
7:  Insert space between digits and letters
8:  Remove all characters except letters, digits, and apostrophes
9:  Collapse multiple spaces into a single space
10: Split text into tokens
11: Remove surrounding quotes from each token
12: Remove stopwords (except "not" and "no")
13: Remove tokens with length less than ar equal 1
14: Apply Porter stemming to remaining tokens
15: Return final token list
\end{lstlisting}
\vspace{1cm}
\Needspace{4\baselineskip}
\noindent\textbf{Transform text using TF-IDF vectorization:}\par
The cleaned text is converted into numerical feature vectors using a TF-IDF vectorizer, which captures the importance of words based on their frequency across documents.
\vspace{1cm}
\Needspace{4\baselineskip}
\noindent\textbf{Encode target labels:}\par
Categorical target labels are converted into numerical form using label encoding. Missing target values are filled using the most frequent class.
\vspace{1cm}
\Needspace{4\baselineskip}
\noindent\textbf{Save preprocessing artifacts:}\par
All fitted preprocessing objects such as the TF-IDF vectorizer and label encoder are saved to disk for reuse during inference.
\vspace{1cm}
\Needspace{4\baselineskip}
\noindent\textbf{Generate preprocessing report (optional):}\par
If enabled, a detailed report is generated containing preprocessing steps, dataset statistics, and visualizations for analysis and reproducibility.
\vspace{1cm}
\Needspace{4\baselineskip}
\noindent\textbf{Return processed datasets:}\par
The final processed training and validation datasets are returned, including transformed features and encoded labels, ready for model training.
\vspace{2cm}
\begin{lstlisting}[style=algostyle,caption={Common Preprocessing Pipeline},label={alg:common_preprocessing}]
Input: Raw dataframe D
Output: X_train, y_train, X_val, y_val

1:  Compute data overview statistics
2:  Remove duplicate rows from dataset
3:  Split dataset into:
        X_train, X_val, y_train, y_val
4:  Generate exploratory graphs (if enabled)
5:  Apply feature preprocessing on X_train and X_val
6:  Apply target preprocessing on y_train and y_val
7:  If is_report = True then
8:     Execute report generation
9: End if
10: Return X_train, y_train, X_val, y_val
\end{lstlisting}
\vspace{1cm}
\Needspace{4\baselineskip}
\noindent\textbf{Testing Preprocessing Steps}\par

\begin{enumerate}
  \item Load saved training metadata from \texttt{data\_info.pkl} and recover the original \texttt{text\_col} and \texttt{target\_col}.
  \item Validate that the test DataFrame contains the text column. If the target column is absent, enable feature-only mode.
  \item Normalize string values by lowercasing the entire DataFrame using \texttt{lower\_data}. 
  \item Fill missing text values with the empty string and transform test text using the saved TF-IDF transformer.
  \item If target labels exist, fill missing labels with the saved training mode and apply the saved label encoder.
  \item Return preprocessed text features and labels (or \texttt{None} for labels when feature-only mode is enabled).
\end{enumerate}
\vspace{1cm}
\Needspace{4\baselineskip}
\noindent\textbf{Constructor Initialization and Artifact Loading:}\par

The testing constructor loads \texttt{data\_info.pkl} from the saved folder path and recovers the original \texttt{text\_col}, \texttt{target\_col}, and \texttt{target\_mode}. It verifies that the text column exists in the test DataFrame and sets \texttt{features\_only=true} if the target column is missing.
\vspace{1cm}
\Needspace{4\baselineskip}
\noindent\textbf{Text Testing Preprocessing Parameters:}\par

\begin{table}[!htbp]
\centering
\caption{Text Testing Preprocessing Parameters}
\label{tab:text_testing_params}
\begin{tabular}{p{0.30\linewidth}p{0.65\linewidth}}
\toprule
\textbf{Parameter} & \textbf{Description} \\
\midrule

\texttt{df} & Input pandas DataFrame containing test examples.\\

\texttt{folder\_path} & Directory path containing saved training artifacts and metadata.\\

\bottomrule
\end{tabular}
\end{table}
\vspace{1cm}
\Needspace{4\baselineskip}
\noindent\textbf{Feature Preprocessing:}\par

The feature preprocessing method fills missing values in the text column with empty strings and applies the loaded TF-IDF transformer from \texttt{training\_preprocessor.pkl}. This yields the final sparse feature matrix for test time.
\vspace{1cm}
\Needspace{4\baselineskip}
\noindent\textbf{Target Preprocessing:}\par

If the test set contains the target column, missing targets are replaced with the saved training mode and then transformed with the loaded label encoder from \texttt{target\_preprocessor.pkl}. If the target column is missing, the method returns \texttt{None} for the labels.
\vspace{1cm}
\Needspace{4\baselineskip}
\noindent\textbf{Common Preprocessing Flow:}\par

The test pipeline first lowercases all string values in the DataFrame, then executes feature preprocessing and target preprocessing sequentially. It returns the transformed text features and encoded labels.
\vspace{0.5cm}

\begin{lstlisting}[style=algostyle,caption={Text Testing Preprocessing Pipeline},label={alg:text_testing_pipeline}]
Input: test DataFrame D, saved folder path
Output: X_test, y_test or None

1:  Load data_info.pkl and recover text_col, target_col, target_mode
2:  Validate that text_col exists in D
3:  If target_col not in D then set features_only = True
4:  Lowercase all string values in D
5:  Fill missing values in D[text_col] with empty string
6:  Apply saved TF-IDF transformer to D[text_col]
7:  If features_only then
8:      y_test = None
9:  else
10:     Fill missing values in D[target_col] with target_mode
11:     Transform targets with the saved label encoder
12: End if
13: Return X_test, y_test
\end{lstlisting}
\vspace{1cm}
\Needspace{4\baselineskip}
\noindent\textbf{Saved artifacts used by text testing:}\par

\begin{table}[!htbp]
\centering
\caption{Artifacts consumed by Text Testing Preprocessing}
\begin{tabular}{p{0.30\linewidth}p{0.65\linewidth}}
\toprule
\textbf{Artifact} & \textbf{Purpose}\\
\midrule
\texttt{data\_info.pkl} & Contains saved metadata: \texttt{text\_col}, \texttt{target\_col}, and \texttt{target\_mode}.\\
\texttt{training\_preprocessor.pkl} & Saved TF-IDF vectorizer used to transform test text.\\
\texttt{target\_preprocessor.pkl} & Saved LabelEncoder used to transform test labels.\\
\bottomrule
\end{tabular}
\end{table}
\subsubsection{Image Dataset}

The image dataset includes classification tasks that operate on folder-structured image data. The classification image preprocessing pipeline is implemented in `ClassificationImageTrainingPreprocessing` and inherits common image preprocessing logic from `ImageTrainingPreprocessing`, which itself inherits from `PreprocessingBase`.

\vspace{1cm}
\Needspace{4\baselineskip}
\noindent\textbf{Classification image training preprocessing parameters:}\par
\begin{table}[!htbp]
\centering
\caption{Classification Image Training Preprocessing Parameters}
\label{tab:classification_image_training_params}
\begin{tabular}{p{0.30\linewidth}p{0.65\linewidth}}
\toprule
\textbf{Parameter} & \textbf{Description} \\
\midrule

\texttt{training\_folder\_images} & Path to the folder containing training image class subdirectories. \\

\texttt{folder\_path} & Directory path used to save report artifacts and metadata. \\

\texttt{validation\_folder\_images} & Optional path to a separate validation folder with the same class subdirectories. \\

\texttt{split\_training} & Boolean flag to split training images into training/validation when no validation folder is provided. \\

\texttt{val\_size} & Validation split fraction when splitting the training data (default: 0.2). \\

\texttt{random\_state} & Random seed for data splitting, shuffling, and report reproducibility (default: 42). \\

\texttt{images\_size} & Output image size tuple, e.g. (224, 224). \\

\texttt{augment} & Enable image augmentation during training data loading. \\

\texttt{batch\_size} & Batch size used by the PyTorch DataLoader (default: 16). \\

\texttt{augmentation\_probability} & Probability of applying augmentation to each training image (default: 0.5). \\

\texttt{horizontal\_flip} & Enable random horizontal flip augmentation. \\

\texttt{vertical\_flip} & Enable random vertical flip augmentation. \\

\texttt{rotation} & Enable random rotation augmentation. \\

\texttt{rotation\_angle} & Maximum rotation angle in degrees (default: 30). \\

\texttt{brightness} & Enable brightness augmentation. \\

\texttt{brightness\_factors} & Brightness factor range tuple (default: (0.8, 1.2)). \\

\texttt{contrast} & Enable contrast augmentation. \\

\texttt{contrast\_factors} & Contrast factor range tuple (default: (0.8, 1.2)). \\

\bottomrule
\end{tabular}
\end{table}

\vspace{1cm}
\Needspace{4\baselineskip}
\noindent\textbf{Image preprocessing validation rules:}\par
The classification image preprocessing pipeline validates all configuration parameters and input folders before proceeding.
\begin{itemize}
  \item `training\_folder\_images` must not be null and must exist.
  \item If `validation\_folder\_images` is provided, it must exist, be different from the training folder, and `split\_training` must be false.
  \item Training folder must contain at least one class subdirectory.
  \item If a validation folder is provided, it must contain at least one class subdirectory.
  \item All validation classes must also exist in the training classes.
  \item `val\_size` must be float in $(0,0.5]$.
  \item `random\_state` must be integer or None.
  \item `images\_size` must be a tuple of two integers between 32 and 1024.
  \item `batch\_size` must be a positive integer.
  \item `augment`, `horizontal\_flip`, `vertical\_flip`, `rotation`, `brightness`, and `contrast` must each be boolean.
  \item `augmentation\_probability` must be between 0 and 1.
  \item `rotation\_angle` must be a non-negative number.
  \item `brightness\_factors` and `contrast\_factors` must each be tuples of two numeric values.
  \item The class-to-label and label-to-class mappings are saved as `class2label\_mapping.pkl` and `label2class\_mapping.pkl`.
  \item Image size metadata is saved as `data\_info.pkl`.
\end{itemize}

\vspace{1cm}
\Needspace{4\baselineskip}
\noindent\textbf{Classification image preprocessing flow:}\par
\begin{enumerate}
  \item Enumerate class subdirectories from the training folder and store the class list.
  \item Build class-label mappings and persist them.
  \item For each image path in training and optional validation folders, validate the file is a supported image.
  \item Collect valid image paths and labels, and record invalid image paths for reporting.
  \item Generate a data overview summary with class counts, sample file paths, and invalid image information.
  \item Split training data into training and validation sets if `split\_training=True`, otherwise use the provided validation folder.
  \item Build PyTorch `DataLoader` objects for training and validation data using `ClassificationImageDataset`.
  \item Apply augmentation to training images when `augment=True` with probability `augmentation\_probability`. The augmentation pipeline may include:
  \begin{itemize}
    \item random horizontal flip
    \item random vertical flip
    \item random rotation within `rotation\_angle`
    \item brightness adjustment using `brightness\_factors`
    \item contrast adjustment using `contrast\_factors`
  \end{itemize}
  \item Generate graphs that describe the prepared dataset and batch behavior:
  \begin{itemize}
    \item label distribution graph showing the count of images per class,
    \item sample images graph showing example images from each class,
    \item example batch graph showing a batch of images after loader creation and augmentation.
  \end{itemize}
  \item Execute `ImagePreprocessingReport` to render the saved report artifacts.
\end{enumerate}

\vspace{1cm}
\Needspace{4\baselineskip}
\noindent\textbf{Saved artifacts for classification image preprocessing:}\par
\begin{itemize}
  \item \texttt{data\_info.pkl}: saved image size metadata.
  \item \texttt{class2label\_mapping.pkl}: mapping from class name to integer label.
  \item \texttt{label2class\_mapping.pkl}: mapping from integer label back to class name.
\end{itemize}

\vspace{1cm}
\Needspace{4\baselineskip}
\noindent\textbf{End-to-end pipeline summary:}\par
The `common\_preprocessing()` method in `ClassificationImageTrainingPreprocessing` performs class mapping, data preparation, overview generation, train/validation splitting, loader creation, graph generation, and finally report execution.

\vspace{1cm}
\Needspace{4\baselineskip}
\subsubsection{Classification Image Testing}
The testing preprocessing pipeline is implemented in `ClassificationImageTestingPreprocessing` and is used after training is complete to prepare new image paths for model inference.

\vspace{1cm}
\Needspace{4\baselineskip}
\noindent\textbf{Classification image testing preprocessing overview:}\par
The testing pipeline accepts a list of image paths and optional class names, validates each image path, loads saved training metadata, and returns a `DataLoader` for inference.

\begin{lstlisting}[style=algostyle,caption={Classification Image Testing Preprocessing Pipeline},label={alg:classification_image_testing_pipeline}]
Input: image paths list P, optional class names list C, saved folder path
Output: testing DataLoader, invalid image paths

1:  Validate P is a non-empty list of image paths
2:  If C is provided, validate it is a list with length equal to P
3:  For each path in P:
4:      If path is a valid image, add to valid image list
5:      Else add to invalid image list
6:  If no valid images remain, raise an error
7:  Load `data_info.pkl` from folder path and recover `image_size`
8:  Load `class2label_mapping.pkl` from folder path
9:  If C is provided then
10:     Convert each provided class name to its integer label
11: End if
12:  Build `ClassificationImageDataset` with `augment=False`
13:  Create a PyTorch `DataLoader` with `shuffle=False` and `collect_function`
14:  Return the testing loader and invalid image paths
\end{lstlisting}

\vspace{1cm}
\Needspace{4\baselineskip}
\noindent\textbf{Testing artifacts consumed:}\par
\begin{itemize}
  \item \texttt{data\_info.pkl}: used to restore `image\_size` for test images.
  \item \texttt{class2label\_mapping.pkl}: used to convert class names to label integers when provided.
\end{itemize}

\vspace{1cm}
\Needspace{4\baselineskip}
\subsubsection{Segmentation Image Training}
The segmentation image training preprocessing class is implemented in `SegmentationImageTrainingPreprocessing` and extends the base image preprocessing pipeline with paired mask handling.

\vspace{1cm}
\Needspace{4\baselineskip}
\noindent\textbf{Segmentation image training preprocessing parameters:}\par
\begin{table}[!htbp]
\centering
\caption{Segmentation Image Training Preprocessing Parameters}
\label{tab:segmentation_image_training_params}
\begin{tabular}{p{0.30\linewidth}p{0.65\linewidth}}
\toprule
\textbf{Parameter} & \textbf{Description} \\
\midrule

\texttt{training\_folder\_images} & Path to the folder containing training images. \\

\texttt{training\_folder\_masks} & Path to the folder containing training masks paired with the training images. \\

\texttt{folder\_path} & Directory path used to save report artifacts and metadata. \\

\texttt{validation\_folder\_images} & Optional path to a separate validation image folder. \\

\texttt{validation\_folder\_masks} & Optional path to the corresponding validation mask folder. \\

\texttt{split\_training} & Boolean flag to split training images into training/validation when no validation folder is provided. \\

\texttt{val\_size} & Validation split fraction when splitting the training data (default: 0.2). \\

\texttt{random\_state} & Random seed for data splitting, shuffling, and report reproducibility (default: 42). \\

\texttt{images\_size} & Output image and mask size tuple, e.g. (224, 224). \\

\texttt{augment} & Enable image augmentation during training data loading. \\

\texttt{batch\_size} & Batch size used by the PyTorch DataLoader (default: 16). \\

\texttt{augmentation\_probability} & Probability of applying augmentation to each training image/mask pair (default: 0.5). \\

\texttt{horizontal\_flip} & Enable random horizontal flip augmentation. \\

\texttt{vertical\_flip} & Enable random vertical flip augmentation. \\

\texttt{rotation} & Enable random rotation augmentation. \\

\texttt{rotation\_angle} & Maximum rotation angle in degrees (default: 30). \\

\texttt{brightness} & Enable brightness augmentation. \\

\texttt{brightness\_factors} & Brightness factor range tuple (default: (0.8, 1.2)). \\

\texttt{contrast} & Enable contrast augmentation. \\

\texttt{contrast\_factors} & Contrast factor range tuple (default: (0.8, 1.2)). \\

\texttt{binary\_mask\_threshold} & Threshold used to binarize mask pixel values (default: 128). \\

\bottomrule
\end{tabular}
\end{table}

\vspace{1cm}
\Needspace{4\baselineskip}
\noindent\textbf{Segmentation preprocessing validation rules:}\par
The segmentation training pipeline validates all folder paths, paired masks, and augmentation parameters before proceeding.
\begin{itemize}
  \item `training\_folder\_images` and `training\_folder\_masks` must not be null, must exist, and must be different folders.
  \item `binary\_mask\_threshold` must be an integer between 0 and 255.
  \item If `validation\_folder\_images` is provided, `validation\_folder\_masks` must also be provided, and both folders must exist and be different.
  \item `training\_folder\_images` and `training\_folder\_masks` must contain matching filenames for valid image/mask pairs.
  \item `val\_size` must be float in $(0,0.5]$.
  \item `random\_state` must be integer or None.
  \item `images\_size` must be a tuple of two integers between 32 and 1024.
  \item `batch\_size` must be a positive integer.
  \item `augment`, `horizontal\_flip`, `vertical\_flip`, `rotation`, `brightness`, and `contrast` must each be boolean.
  \item `augmentation\_probability` must be between 0 and 1.
  \item `rotation\_angle` must be a non-negative number.
  \item `brightness\_factors` and `contrast\_factors` must each be tuples of two numeric values.
\end{itemize}

\vspace{1cm}
\Needspace{4\baselineskip}
\noindent\textbf{Segmentation image preprocessing flow:}\par
\begin{enumerate}
  \item Validate training image and mask folder paths and optional validation folder pairs.
  \item Verify mask folders are distinct from the corresponding image folders.
  \item Save `data\_info.pkl` containing `image\_size` and `binary\_mask\_threshold`.
  \item Match image and mask filenames in each folder pair, keeping only valid pairs and recording invalid image or mask paths.
  \item Generate an overview summary of training and optional validation data, including sample file pairs and invalid path counts.
  \item Split the training pairs into training and validation sets if `split\_training=True`, otherwise use the provided validation folder pairs.
  \item Build PyTorch `DataLoader` objects with `SegmentationImageDataset` for training and validation; training loader uses augmentation while validation loader does not.
  \item Apply augmentation to training image/mask pairs when `augment=True` with probability `augmentation\_probability`. The augmentation pipeline may include:
  \begin{itemize}
    \item random horizontal flip,
    \item random vertical flip,
    \item random rotation within `rotation\_angle`,
    \item brightness adjustment using `brightness\_factors`,
    \item contrast adjustment using `contrast\_factors`.
  \end{itemize}
  \item Generate graphs that describe dataset structure and loader batches, including folder summaries, sampled image/mask pairs, and loader batch examples.
  \item Execute `ImagePreprocessingReport` to render the segmentation training report artifacts.
\end{enumerate}

\vspace{1cm}
\Needspace{4\baselineskip}
\noindent\textbf{Saved artifacts for segmentation image preprocessing:}\par
\begin{itemize}
  \item \texttt{data\_info.pkl}: saved image size metadata and `binary\_mask\_threshold`.
\end{itemize}

\vspace{1cm}
\Needspace{4\baselineskip}
\subsubsection{Segmentation Image Testing}
The testing preprocessing pipeline is implemented in `SegmentationImageTestingPreprocessing` and is used after training is complete to prepare new image and mask paths for inference or evaluation.

\vspace{1cm}
\Needspace{4\baselineskip}
\noindent\textbf{Segmentation image testing preprocessing overview:}\par
The testing pipeline accepts a list of image paths and optional mask paths, validates each path pair, loads saved training metadata, and returns a `DataLoader` for inference.

\begin{lstlisting}[style=algostyle,caption={Segmentation Image Testing Preprocessing Pipeline},label={alg:segmentation_image_testing_pipeline}]
Input: image paths list P, optional mask paths list M, saved folder path
Output: testing DataLoader, invalid image paths

1:  Validate P is a non-empty list of image paths
2:  If M is provided, validate it is a list with length equal to P
3:  For each path in P:
4:      If path is a valid image and corresponding mask path is valid (when provided), add to valid pairs
5:      Else add the image path to invalid image list
6:  If no valid images remain, raise an error
7:  Load `data_info.pkl` from folder path and recover `image_size` and `binary\_mask\_threshold`
8:  Build `SegmentationImageDataset` with `augment=False`
9:  Create a PyTorch `DataLoader` with `shuffle=False` and `collect_function_segmentation`
10: Return the testing loader and invalid image paths
\end{lstlisting}

\vspace{1cm}
\Needspace{4\baselineskip}
\noindent\textbf{Segmentation testing artifacts consumed:}\par
\begin{itemize}
  \item \texttt{data\_info.pkl}: used to restore `image\_size` and `binary\_mask\_threshold` for test images.
\end{itemize}
